\documentclass[12pt,a4paper]{article}
\usepackage[utf8]{inputenc}
\usepackage[T1]{fontenc}
\usepackage{geometry}
\usepackage{hyperref}
\usepackage{enumitem}
\usepackage{booktabs}
\usepackage{float}

\geometry{margin=1in}

\title{AI Storytelling Adventure Game\\
\large System Testing and Evaluation Report}
\author{Dumitru Nirca}
\date{28th February 2026}

\begin{document}

\maketitle

\section{Overview}

Testing covered functional correctness, integration, performance, output quality, and user experience. Evaluation relied exclusively on AI agents (GPT-4, Claude 3.5 Sonnet, Gemini); no human evaluators were involved. Each criterion was rated on a 1--5 scale, where 1 is lowest and 5 is highest; all reported scores are out of 5. The report first assesses whether the system works as a whole (pass/fail), then presents the request lifecycle, evaluation criteria, performance benchmarks, AI agent methodology and results, and narrative content examples from the gameplay sessions that produced the evaluation samples.

\section{System-Wide Assessment: Pass/Fail}

Before criterion-based evaluation, the system was assessed as a whole across multiple dimensions to establish whether it functions correctly.

\textbf{Functional testing}: All three services (ComfyUI, Ollama, Flask) start and communicate as required. Startup sequence verified: ComfyUI loads checkpoint (15--20 seconds), Ollama loads on first request, Flask serves the interface. Pre-game selection (start place, character, companion) works; full play loop completes: player submits choice or free text, receives narrative and choices, text types first then image appears (or graceful failure if ComfyUI unavailable). Ollama unavailability displays a dedicated error page.

\textbf{Integration testing}: Flask--Ollama calls succeed; Flask--ComfyUI API (workflow submission, polling, file retrieval) succeeds; client polling and image display succeed. Session state persists across turns; conversation history is correctly reconstructed and passed to the LLM.

\textbf{Error handling}: Unavailable Ollama returns error JSON and client message; unavailable ComfyUI stops polling and continues without image. No crashes observed under expected failure conditions.

\textbf{Output quality}: Text generation produces coherent, kid-friendly scene descriptions and numbered choices; image generation produces relevant children's storybook imagery; story-led prompt enhancement matches scene content reliably for the first 3--5 turns. Kid-friendly safeguards (vocabulary constraints, ``never scary or violent,'' storybook tone, image style suffix, negative prompts) were verified during gameplay sessions; no inappropriate content was observed within the tested scope.

\textbf{Verdict}: \textbf{PASS}. The system functions end-to-end. All components integrate correctly, the full game loop completes, and degradation (narrative drift, image mismatches) is predictable and confined to known architectural limits.

\section{Testing Dimensions}

Testing touched the following areas:

\textbf{Functional}: Service startup, API availability, request/response correctness.

\textbf{Integration}: Flask--Ollama, Flask--ComfyUI, client--server polling.

\textbf{Performance}: Text generation (2--4 seconds), image generation (28 seconds), text-first-then-image display order with polling.

\textbf{Output quality}: Narrative coherence, image relevance, choice validity.

\textbf{UX}: Typing animation, image fade-in, choice buttons, free-text input.

\section{Request Lifecycle}

Understanding the full lifecycle helps clarify what was tested. Before the story begins, the player completes a three-step selection (start place, character, companion) via \texttt{GET /} with query parameters. When all three are chosen, the session is created, the first scene is generated by Ollama, and the page loads immediately; the first-scene image is generated in a background thread and the client polls \texttt{/get\_image/scene\_initial} until ready. Text displays first (typing animation), then the image when available.

When a player clicks a choice button or submits free text (during play), the following sequence occurs:

\begin{enumerate}
    \item The browser sends a POST request to \texttt{/make\_choice} with a JSON body containing either \texttt{choice} (an integer) or \texttt{free\_text} (a string).
    \item Flask retrieves the player's session, reconstructs the conversation history, and appends the player's input as a user message.
    \item The updated conversation history is sent to the Ollama \texttt{/api/chat} endpoint. This call is synchronous and blocks the Flask thread until Ollama responds, typically 2--4 seconds.
    \item The Ollama response is parsed: the scene text is extracted by stripping markdown artifacts, and the numbered choices are extracted by scanning for lines matching the pattern \texttt{N. option text}.
    \item A new \texttt{scene\_id} is assigned (e.g.\ \texttt{scene\_7}) and added to \texttt{pending\_images} with \texttt{None} to indicate generation in progress.
    \item A daemon thread calls the ComfyUI API to queue image generation and blocks until completion.
    \item Flask returns a JSON response to the browser containing scene text, choices, player status, and scene ID (before the image thread has made progress).
    \item The browser starts the typing animation (45ms per character) and polls \texttt{/get\_image/scene\_7} every 200ms.
    \item The image thread completes after approximately 28 seconds, moves the file to \texttt{static/images/}, and updates \texttt{pending\_images}.
    \item On the next poll, the image URL is returned; the browser inserts the image with a CSS fade-in.
\end{enumerate}

Error handling at steps 3 and 6 checks service reachability. If Ollama is unavailable, the server returns error JSON and the client displays a ``service unavailable'' message. If ComfyUI is unavailable, \texttt{pending\_images[scene\_id]} is set to \texttt{None} permanently and the poll endpoint returns a failure flag; the client stops polling and continues without an image.

\section{Evaluation Criteria}

Evaluation was conducted using four criteria. Each criterion was rated on a 1--5 Likert scale, where 1 represents the lowest performance and 5 the highest; all reported scores (e.g., 3.9) are therefore out of 5.

\textbf{Enjoyment} (including user engagement): Whether the overall experience is engaging and worth continuing, and whether the interface and mechanics (typing animation, audio, choice presentation) keep the player invested.

\textbf{Narrative Coherence}: Whether the AI-generated story is consistent within and across turns, with logical cause and effect and stable characters.

\textbf{Image Quality}: Whether the generated images are visually acceptable and relevant to the scene described in the text.

\textbf{System Responsiveness} (including synchronisation): How quickly the system responds to player input, including text generation latency and delays before the image appears. The system displays text first (typing animation), then the image when ready.

\section{Performance Metrics and Benchmarks}

Table~\ref{tab:performance_metrics} presents comprehensive performance metrics collected during system testing.

\begin{table}[H]
\centering
\caption{System Performance Metrics and Benchmarks}
\label{tab:performance_metrics}
\begin{tabular}{lc}
\toprule
\textbf{Metric} & \textbf{Value} \\
\midrule
Average Text Generation Time & 2--4 seconds \\
Average Image Generation Time & 28 seconds \\
Text Typing Animation Duration (600--700 chars) & 27--31 seconds \\
Display Order & Text first, then image \\
Image Polling Interval & 200ms \\
Cache Hit Rate (repeated scenes) & $\sim$15--20\% \\
Average Scene Load Time (first scene) & 30--35 seconds \\
Average Scene Load Time (subsequent scenes) & 30--33 seconds \\
Conversation History Window & 6 messages \\
Maximum Text Length Generated & 700 characters \\
Minimum Text Length Generated & 450 characters \\
Average Text Length Generated & 625 characters \\
Image Generation Queue Delay (peak) & 5--8 seconds \\
\bottomrule
\end{tabular}
\end{table}

The bottleneck is image generation at 28 seconds; text generation at 2--4 seconds is effectively instantaneous from the user's perspective. The system displays text first (typing animation), then the image when ready. Shorter scenes (e.g.\ 450 characters, ~20 seconds display time) may finish typing with a wait before the image arrives; longer scenes (600--700 characters, ~27--31 seconds) typically have the image ready by the time typing completes. The fade-in transition smooths the appearance.

Synchronisation calibration was tuned during testing: at 30ms per character, a 600-character response displays in approximately 18 seconds --- a systematic 10-second gap before the image arrives. Increasing to 45ms raised display time to approximately 29 seconds, matching image generation. Residual variance of approximately 4 seconds arises from the model generating a variable number of characters despite the 600--700 character instruction (standard deviation $\approx$85 characters across 50 sampled responses).

Cache hits at 15--20\% arise primarily when similar player actions produce identical story-led prompts. The hash is computed on the enhanced prompt rather than the raw text, so two different narrative phrasings that produce the same enhanced prompt share a cache entry. Queue delays of 5--8 seconds occur when a second scene is requested before the first image has finished generating; the second request waits in ComfyUI's queue.

\section{AI Agent Testing Methodology}

Evaluation was conducted solely by AI agents; no human participants or evaluators were involved. Representative text and image samples were collected from five gameplay sessions (15--25 minutes each), covering village and forest exploration, castle and seaside settings, companion-focused sequences, and a free-form session using custom text input. These samples were provided to GPT-4, Claude 3.5 Sonnet, and Gemini, each asked to rate the system against the same four criteria on a 1--5 Likert scale (1 = lowest, 5 = highest; all scores reported as out of 5) with written justifications.

Using three models reduces the risk of any single model's systematic tendencies dominating the results. The AI agents evaluated the same materials under the same criteria, enabling cross-model agreement analysis.

\section{AI Agent Testing Results}

Table~\ref{tab:ai_testing_metrics} summarises the ratings from each AI agent across the four criteria. All scores are out of 5 (1 = lowest, 5 = highest).

\begin{table}[H]
\centering
\caption{AI Agent Testing Metrics Summary (scores out of 5)}
\label{tab:ai_testing_metrics}
\begin{tabular}{lcccc}
\toprule
\textbf{Evaluation Criterion} & \textbf{GPT-4} & \textbf{Claude 3.5} & \textbf{Gemini} & \textbf{Average} \\
\midrule
Enjoyment (incl.\ user engagement) & 4.2 & 3.9 & 4.0 & 4.0 \\
Narrative Coherence & 3.7 & 3.6 & 3.8 & 3.7 \\
Image Quality & 4.3 & 4.0 & 4.0 & 4.1 \\
System Responsiveness (incl.\ synchronisation) & 3.8 & 3.6 & 3.7 & 3.7 \\
\midrule
\textbf{Overall Average} & \textbf{4.0} & \textbf{3.8} & \textbf{3.9} & \textbf{3.9/5} \\
\bottomrule
\end{tabular}
\end{table}

All three models agreed on the primary strength and primary weakness. Strengths: image quality (4.0--4.3 out of 5) and enjoyment including interface engagement (3.9--4.2 out of 5); the retro interface and typing animation were consistently noted as positive contributors. Weakness: narrative coherence degradation in longer sessions, with Claude 3.5 specifically identifying ``character motivations shifted unexpectedly'' and ``plot threads were introduced but not resolved.'' All three models noted that the text-first-then-image display order provides a clear narrative flow; subsequent scenes may experience delays when ComfyUI queue processing backs up. Gemini noted that the 6-message context limit may be insufficient for complex narratives.

Because evaluation relied exclusively on AI agents, the detailed written justifications for each score provided specific diagnostic information (e.g., naming particular narrative inconsistencies) that ratings alone would not have captured.

\section{Narrative Content Examples}

To give a concrete sense of the AI's output quality, the following examples were drawn from the gameplay sessions whose text and image samples were provided to the AI agents for evaluation.

\textbf{Town exploration}: The initial scene for a town session read: \textit{``You arrive in the bustling market square of Thornhaven as the sun climbs past midday. The smell of fresh bread mingles with sawdust from a nearby carpenter's workshop. A town crier is reading a notice beside the well, and a group of merchants are arguing loudly near the eastern gate.''} This was accompanied by an image of a medieval market square with warm afternoon lighting, correctly identified via the ``town'' keyword. The three generated choices were: (1) approach the town crier, (2) speak with the merchants, (3) look for a place to rest. The scene was coherent and the image matched.

\textbf{Dungeon session, early}: After choosing to descend into a dungeon, the model produced: \textit{``The air grows colder as you descend the stone steps. Torches in rusted brackets cast uneven orange light across the corridor walls, and you can hear the distant sound of dripping water. A wooden door stands ajar to your left, and the corridor continues into shadow ahead.''} The image showed a dark stone corridor with torch sconces, matching the scene well.

\textbf{Dungeon session, sixth turn}: By the sixth interaction in the same session, some drift was evident. The model described a room with a ``great stained glass window'' inside a dungeon, which is inconsistent with the underground setting established earlier. The image generated for this scene showed an outdoor cathedral exterior, which had been triggered by the word ``window'' being processed as an architectural keyword. This is a clear example of the context window limitation and the limitations of keyword-based image prompting.

\textbf{Free-text input example}: After the standard choices were presented, the player typed: ``I try to pick the lock on the iron chest.'' The model responded with a scene describing the lockpicking attempt and introduced a lock quality variable (``the lock is old but sturdy''), which it had not established previously. The response was plausible but introduced game mechanics (skill checks, lock quality) that were inconsistent with the narrative's established tone. This illustrates the free-text coherence issue noted in the limitations.

These examples show that the system works well for the first 3--5 interactions and degrades in specific ways that are predictable from the architecture.

\end{document}
